\section{Bounded Booleanization and finite search}
\label{sec:booleanization}

\subsection{Boolean predicates}

A numerical predicate \(P(\bar x)\) is represented by a primitive-recursive
function \(b_P(\bar x)\in\B\).  We write
\[
  [P](\bar x) \quad\text{for the equation}\quad b_P(\bar x)=1.
\]
The Boolean operations are fixed primitive-recursive functions satisfying the
truth tables for \(0\) and \(1\).  Their correctness is proved by the two
numerical cases.  Equality and order are also represented by Boolean functions;
for example, \(x\le y\) can be tested by whether truncated subtraction
\(x\dotminus y\) is zero.

A \emph{bounded formula} is built from atomic equations between
primitive-recursive terms, propositional connectives, and quantifiers
\(\forall i<B(\bar x)\) and \(\exists i<B(\bar x)\).  The bound and all terms
are fixed primitive-recursive functions.

\subsection{Finite folds}

For a Boolean predicate \(b(i,\bar x)\), define
\begin{align*}
 \All_b(0,\bar x)&=1,\\
 \All_b(n+1,\bar x)&=\All_b(n,\bar x)\wedge b(n,\bar x),\\[3pt]
 \Any_b(0,\bar x)&=0,\\
 \Any_b(n+1,\bar x)&=\Any_b(n,\bar x)\vee b(n,\bar x).
\end{align*}
Both are primitive recursive.  Define a first-witness search
\(\Search_b(n,\bar x)\) that returns the least \(i<n\) with
\(b(i,\bar x)=1\), and returns \(n\) when no such witness exists.  It can be
computed by a primitive recursion carrying the current candidate and a found
bit.

\begin{theorem}[Correctness of Boolean folds]
\label{thm:boolean-fold}
For every fixed Boolean predicate \(b\), iPRA proves:
\begin{align*}
 \All_b(n,\bar x)=1
  &\leftrightarrow \forall i\,(i<n\to b(i,\bar x)=1),\\
 \Any_b(n,\bar x)=1
  &\leftrightarrow \exists i\,(i<n\wedge b(i,\bar x)=1),
\end{align*}
and, under the existential premise,
\[
 \Search_b(n,\bar x)<n,
 \qquad b(\Search_b(n,\bar x),\bar x)=1,
\]
with no smaller witness.
\end{theorem}

\begin{proof}
We spell out the proof because it illustrates the permitted use of
quantifier-free induction.

For the forward direction of the universal fold, keep \(i\) free and induct on
\(n\) with the quantifier-free formula
\[
  i<n\wedge \All_b(n,\bar x)=1\;\to\;b(i,\bar x)=1.
\]
The relation \(i<n\) and the Boolean equation are atomic equations between
primitive-recursive terms.  At zero the antecedent is false.  At a successor,
if \(i<n+1\), numerical case analysis gives either \(i<n\), where the induction
hypothesis applies to the prefix fold, or \(i=n\), where the last Boolean
factor gives \(b(n,\bar x)=1\).  Universal introduction over \(i\) is performed
after the induction.

For the reverse direction, place
\(H:=\forall i(i<n\to b(i,\bar x)=1)\) in the context and induct on the fold
length \(m\le n\) with the quantifier-free formula
\(\All_b(m,\bar x)=1\).  The step specializes \(H\) to \(m\), uses the
recursion equation of \(\All_b\), and the Boolean truth table for conjunction.
Equivalently, one may prove the open theorem with \(n\) as induction variable,
keeping \(\forall i(i<n\to\cdots)\) outside the induction predicate and using
its successor decomposition in the context.

The existential fold is dual.  For the forward direction, induction on \(n\)
examines whether the last disjunct is true; the result is either the last index
or a witness obtained from the prefix induction hypothesis.  For the reverse
direction, a supplied witness \(i<n\) is propagated through the fold by
quantifier-free induction on the difference between \(i+1\) and \(n\).

For first-witness search, the invariant states that either the found bit is
zero and every earlier predicate value is zero, or the found bit is one, the
stored index is below the current length, its predicate value is one, and all
smaller indices fail.  Every clause is a Boolean equation or bounded
comparison, hence quantifier free.  At termination, an existential premise
excludes the no-witness branch.  The minimality clause is read directly from
the invariant.
\end{proof}

\subsection{Booleanization theorem}

\begin{theorem}[Bounded Booleanization]
\label{thm:booleanization}
For every fixed bounded formula \(\theta(\bar x)\), one can construct a
primitive-recursive Boolean function \(b_\theta\) such that iPRA proves
\[
  b_\theta(\bar x)=1\quad\leftrightarrow\quad\theta(\bar x).
\]
The construction is effective in the syntax of \(\theta\).
\end{theorem}

\begin{proof}
Proceed by structural recursion on \(\theta\).  Atomic equations use the
Boolean equality test.  Conjunction, disjunction, and implication use the
corresponding Boolean truth functions; intuitionistic logic proves the finite
truth-table equivalences because each input is known to equal zero or one.
Negation is implication to the designated false equation.

For \(\forall i<B(\bar x)\,\psi(i,\bar x)\), apply the induction hypothesis to
obtain \(b_\psi\), and define
\[
 b_\theta(\bar x)=\All_{b_\psi}(B(\bar x),\bar x).
\]
The result follows from \cref{thm:boolean-fold}.  The existential case uses
\(\Any\).  All constructions use composition and primitive recursion, so the
resulting function is represented by a fixed \(\PR\) symbol.  The proof is
uniformly assembled from the component equivalences by substitution and
natural-deduction rules.
\end{proof}

\subsection{Bounded minimization}

Let \(c(i,\bar x)\) be an integer or rational cost with a decidable comparison,
and let \(C(\bar x)>0\) be a finite candidate bound.  Define
\(\operatorname{argmin}_{i<C}c(i,\bar x)\) by primitive recursion, retaining
the least index among equal minima.

\begin{lemma}[Finite minimization]
\label{lem:finite-min}
I P R A proves that the returned index is below \(C\) and has cost no larger
than every candidate below \(C\).  If a Boolean admissibility predicate is
included and at least one candidate is admissible, the same holds among
admissible candidates.
\end{lemma}

\begin{proof}
The accumulator stores the current best index, a found bit, and optionally the
current best cost.  The quantifier-free invariant states that the best index is
admissible when the found bit is one and that every previously inspected
admissible candidate has cost at least the stored cost; tie-breaking adds that
no smaller equal-cost candidate was skipped.  The successor step is one
comparison and a Boolean case split.  The existential admissibility premise
eliminates the no-candidate final state.
\end{proof}

This lemma is the computational core of finite-net residual minimization,
finite maximum and minimum, canonical certificate selection, and the backward
viability algorithms below.

\subsection{Bounded descent}

A common proof has a measure \(m(s)\in\N\) that strictly decreases whenever a
state is not terminal.  If the initial measure is bounded by \(M(u)\), then the
state must terminate within \(M(u)\) steps.

\begin{lemma}[Primitive-recursive descent]
\label{lem:descent}
Suppose fixed primitive-recursive functions \(\mathsf{step}\),
\(\mathsf{terminal}\), and \(m\) satisfy, in iPRA,
\[
 \mathsf{terminal}(s)=0\to m(\mathsf{step}(s))<m(s).
\]
Then the bounded run of length \(m(s)+1\) reaches a terminal state.
\end{lemma}

\begin{proof}
Induct on the step index with the quantifier-free invariant
\[
 \mathsf{terminal}(s_i)=1
 \quad\vee\quad m(s_i)+i\le m(s_0).
\]
If the state is already terminal, keep it fixed.  Otherwise the strict measure
decrease proves the second disjunct at the next stage.  At stage
\(m(s_0)+1\), the nonterminal disjunct would imply a negative natural measure,
so only terminality remains.
\end{proof}

The point is not that all mathematical descent arguments are easy, but that
once a numerical measure and local decrease theorem are supplied, termination
is a routine quantifier-free induction.
